\documentclass[11pt]{article}

\usepackage{amsmath,amssymb,amsfonts}
\usepackage[margin=1in]{geometry}
\usepackage{hyperref}
\usepackage{enumitem}
\usepackage{bm}
\usepackage{booktabs}

\hypersetup{colorlinks=true, linkcolor=blue, citecolor=blue, urlcolor=blue}

\newcommand{\vw}{\bm{v}}
\newcommand{\ew}{\bm{e}}
\newcommand{\KL}{D_{\mathrm{KL}}}
\newcommand{\PMI}{\mathrm{PMI}}
\newcommand{\tr}{\mathrm{tr}}

\title{Word Embeddings from Co-occurrence Statistics:\\A Variational and Statistical-Mechanical Perspective}
\author{}
\date{}

\begin{document}
\maketitle

\begin{abstract}
We present a derivation of word embedding geometry from first principles, starting from a Boltzmann model of word co-occurrence statistics. By modeling the joint probability of a word sequence as an energy-based distribution over count fluctuations, and minimizing the KL divergence between this model and a parameterized softmax family, we show that the optimal embedding vectors satisfy an eigenvalue-like constraint whose solution recovers the well-known result: the dot product of embedding vectors encodes the pointwise mutual information (PMI) between words. We connect this derivation to the established results of Levy and Goldberg~\cite{levy2014}, Pennington et al.~\cite{pennington2014}, and Arora et al.~\cite{arora2016}, providing a unified statistical-mechanical perspective on why word embedding methods work.
\end{abstract}

\section{Introduction}\label{sec:intro}

Word embedding methods such as Word2Vec~\cite{mikolov2013} and GloVe~\cite{pennington2014} map words from a discrete vocabulary of size $N$ into continuous vectors $\vw_i \in \mathbb{R}^d$, with the remarkable property that semantic relationships are encoded in the vector geometry: words that co-occur frequently have large dot products $\vw_i \cdot \vw_j$, and linear directions in embedding space capture analogical structure.

A natural question is: \emph{why} does this work? The answer, developed across several lines of research, is that these algorithms implicitly factorize a matrix of pointwise mutual information (PMI) computed from co-occurrence counts. Levy and Goldberg~\cite{levy2014} showed this explicitly for Word2Vec's skip-gram with negative sampling (SGNS). Arora et al.~\cite{arora2016} provided a generative justification via a latent random-walk model. Hashimoto et al.~\cite{hashimoto2016} framed embeddings as metric recovery in a semantic space.

In this note, we present a self-contained derivation that approaches the problem from a statistical-mechanical perspective. We model word co-occurrence as a Boltzmann distribution whose energy is a quadratic function of count fluctuations (Section~\ref{sec:energy}), perform variational inference by minimizing the KL divergence against a softmax model (Section~\ref{sec:variational}), and show that the stationary-point conditions force the embedding dot products to equal a normalized covariance---i.e., the PMI (Section~\ref{sec:geometry}). We then connect this to the established literature (Section~\ref{sec:connections}), present the precise assumptions and main theorem of the RAND-WALK generative model of Arora et al.~\cite{arora2016} (Section~\ref{sec:randwalk}), and discuss the broader implications (Section~\ref{sec:discussion}).

\section{The Bag-of-Words Energy Model}\label{sec:energy}

\subsection{Setup and notation}

Consider a vocabulary of $N$ tokens indexed by $i = 1, 2, \ldots, N$ and a text window of length $L$ with positions $p = 1, 2, \ldots, L$. Let $y_p \in \{1, \ldots, N\}$ denote the token at position $p$. The \emph{count} of token $i$ in the window is
\begin{equation}\label{eq:count}
    n_i = \sum_{p=1}^{L} \delta_{i, y_p} - P_i\,,
\end{equation}
where $P_i = \langle n_i \rangle$ is the expected count (itself a random variable determined by the underlying distribution), so that $n_i$ measures the \emph{fluctuation} of the count of word $i$ away from its expectation.

\subsection{Boltzmann distribution over co-occurrences}

We model the joint probability of a word sequence as a Boltzmann distribution whose energy is a function of the count fluctuations:
\begin{equation}\label{eq:boltzmann}
    p(w_1, w_2, \ldots, w_L) \propto e^{-\mathcal{E}[N(w_1, w_2, \ldots, w_L)]}\,.
\end{equation}
For the pairwise interaction model, the energy takes the quadratic form
\begin{equation}\label{eq:pairwise}
    \mathcal{E} = \sum_{i,j} G_{ij}\, n_i\, n_j\,,
\end{equation}
so that the pairwise co-occurrence probability factorizes as
\begin{equation}\label{eq:pair-prob}
    p(t_i, t_j) = e^{-G_{ij}\, n_i\, n_j}\,.
\end{equation}
The coupling matrix $G_{ij}$ encodes the strength of interaction between tokens $i$ and $j$. This is the direct analogue of a spin-glass Hamiltonian in statistical mechanics, with the word counts playing the role of spin variables and $G_{ij}$ playing the role of the exchange coupling.

\subsection{Covariance structure and the count fluctuation hyperplane}

The covariance matrix of the embedding vectors defines the \emph{count fluctuation structure}:
\begin{equation}\label{eq:covariance}
    \Sigma = E\, E^\top\,,
\end{equation}
where $E$ is the embedding matrix whose rows are the word vectors. The constraint surface on which the count fluctuations $\{n_i\}$ satisfy the normalization $\sum_i n_i = 0$ (since total count deviations must sum to zero in a fixed-length window) defines a hyperplane in the $N$-dimensional count space---the \emph{count fluctuation hyperplane}. The geometry of word embeddings is determined by the projection of the covariance structure onto this hyperplane.

\section{Variational Inference via KL Divergence}\label{sec:variational}

\subsection{The softmax proposal distribution}

We approximate the true Boltzmann distribution~\eqref{eq:boltzmann} with a tractable softmax model. For each position $p$, the proposal assigns:
\begin{equation}\label{eq:proposal}
    \hat{p}(n_1, n_2, \ldots) = \prod_{p=1}^{L} \frac{e^{\varphi(y_p)}}{\hat{Z}(\varphi)}\,,
\end{equation}
where $\varphi(y_p)$ is a potential assigned to the token at position $p$ and the partition function is
\begin{equation}\label{eq:partition}
    \hat{Z}(\varphi) = \sum_{c=1}^{N} e^{\varphi \cdot \vw_c}\,.
\end{equation}
The goal is to choose the embedding vectors $\{\vw_c\}$ and potentials $\varphi$ such that $\hat{p}$ is as close as possible to the true distribution $p$, in the sense of minimizing the KL divergence $\KL(\hat{p} \,\|\, p)$ from the proposal to the true distribution.

\subsection{KL divergence expansion}

The KL divergence between the proposal and the Boltzmann model is
\begin{equation}\label{eq:kl}
    \KL(\hat{p} \,\|\, p) = \sum_{\{t_i, t_j\}} e^{-G_{ij}\, n_i\, n_j} \log \left( \frac{e^{-G_{ij}\, n_i\, n_j}}{\prod_{p=1}^{L} \hat{Z}(\varphi_p)} \right).
\end{equation}
Expanding the logarithm and collecting terms yields
\begin{equation}\label{eq:kl-expanded}
    \KL = \sum_{i,j} e^{-G_{ij}\, n_i\, n_j} \sum_{p=1}^{N} \hat{n}_p \Big( \log \hat{Z}(\varphi_p) - \log \psi_p \cdot \vw_p \Big)\,,
\end{equation}
where $\hat{n}_p$ are the expected counts under the proposal distribution.

In the softmax formulation, the KL divergence takes the form
\begin{equation}\label{eq:kl-softmax}
    \KL = \sum_{i} P_i \left( \log \frac{\hat{Z}(\vw_i)}{e^{\vw_i \cdot \vw_i}} - \vw_i \cdot \vw_i \right) + \text{const}\,,
\end{equation}
where the sum runs over the vocabulary and $P_i$ are the unigram probabilities.

\subsection{Gradient conditions}\label{sec:gradient}

Taking the derivative of $\KL$ with respect to the embedding vector $\vw_k$ and setting it to zero yields:
\begin{equation}\label{eq:gradient}
    \frac{\partial \KL}{\partial \vw_k} = \sum_{i} e^{-G_{ij}\, n_i\, n_j} \sum_{p} \hat{n}_p \left( \delta_{kp} \frac{\sum_m \vw_m}{\hat{Z}} - \frac{2\vw_k}{\vw_k \cdot \vw_k} \right) - \hat{n}_k = 0\,.
\end{equation}
Under the approximation $\hat{Z}(\varphi_p) \approx \sum_c \vw_c \cdot \vw_c = W_{nk}$, this simplifies to the condition
\begin{equation}\label{eq:stationary}
    -2\hat{n}_k\, \vw_k + \sum_{i} \hat{n}_i\, \vw_k = 0\,.
\end{equation}
This is an eigenvalue-like equation: the embedding vector $\vw_k$ is constrained by a balance between its own count weight and the aggregate count-weighted embedding direction.

\subsection{The matrix constraint}\label{sec:matrix}

Writing the stationary condition~\eqref{eq:stationary} for each $k = 1, 2, \ldots, N$, we obtain a system of linear constraints on the embedding matrix $W$ (whose rows are $\vw_1, \vw_2, \ldots, \vw_N$):
\begin{equation}\label{eq:matrix-constraint}
\begin{aligned}
    k{=}1: &\quad (-n_1,\; n_2,\; n_3,\; \ldots,\; n_N) \begin{pmatrix} \text{---}\, \vw_1 \,\text{---} \\ \text{---}\, \vw_2 \,\text{---} \\ \vdots \\ \text{---}\, \vw_N \,\text{---} \end{pmatrix} = \bm{0}\,, \\[6pt]
    k{=}2: &\quad (n_1,\; -n_2,\; n_3,\; \ldots,\; n_N) \begin{pmatrix} \text{---}\, \vw_1 \,\text{---} \\ \text{---}\, \vw_2 \,\text{---} \\ \vdots \\ \text{---}\, \vw_N \,\text{---} \end{pmatrix} = \bm{0}\,,
\end{aligned}
\end{equation}
and so on for each $k$. Each row of this system singles out the $k$-th word count with a sign flip, requiring that the count-weighted projection onto the embedding space vanishes. This imposes a geometric constraint: the word vectors must be arranged so that count fluctuations are orthogonal to certain directions in embedding space.

\section{Embedding Geometry Encodes Normalized Covariance}\label{sec:geometry}

The central result emerges from the stationary-point analysis. The dot product of embedding vectors satisfies
\begin{equation}\label{eq:main-result}
    \boxed{\;\ew_i \cdot \ew_j = \frac{\Sigma_{ij}}{L\, P_i\, P_j}\;\,}
\end{equation}
where $\Sigma_{ij} = \langle n_i\, n_j \rangle - \langle n_i \rangle \langle n_j \rangle$ is the covariance of the count fluctuations, $L$ is the window length, and $P_i$, $P_j$ are the unigram probabilities. The dot product encodes the \emph{normalized covariance}---a linear measure of the association between tokens $i$ and $j$, measuring how much their co-occurrence deviates from independence.

This is related to, but distinct from, the pointwise mutual information $\PMI(i,j) = \log[p(i,j)/(p(i)p(j))]$, which is a logarithmic measure of association. The two coincide in the regime of weak correlations (where $\log(1+x) \approx x$), but differ in general. The Arora et al.\ RAND-WALK model~\cite{arora2016} yields $\ew_i \cdot \ew_j \propto \PMI(i,j)$ under its own set of assumptions (Section~\ref{sec:randwalk}); the present variational derivation from the bag-of-words energy model gives the normalized covariance directly.

\section{Connection to Established Results}\label{sec:connections}

\subsection{Levy and Goldberg (2014): SGNS as PMI factorization}

Levy and Goldberg~\cite{levy2014} proved that Word2Vec's skip-gram with negative sampling (SGNS) implicitly factorizes the \emph{shifted PMI matrix}. At convergence, the word and context vectors satisfy
\begin{equation}\label{eq:levy-goldberg}
    \vw_i \cdot \bm{c}_j = \PMI(w_i, c_j) - \log k\,,
\end{equation}
where $k$ is the number of negative samples. The proof proceeds by treating each $(w, c)$ entry of the SGNS objective as an independent logistic regression, whose closed-form optimum yields the PMI value. Our variational derivation arrives at the same conclusion from the energy-model side: the KL-optimal embedding dot products encode the co-occurrence structure, which is PMI.

\subsection{GloVe: biases absorb marginals}

The GloVe objective~\cite{pennington2014} fits
\begin{equation}\label{eq:glove}
    \vw_i \cdot \vw_j + b_i + b_j = \log X_{ij}\,,
\end{equation}
where $X_{ij}$ is the co-occurrence count. Since $\log X_{ij} = \PMI(i,j) + \log X_{i\cdot} + \log X_{\cdot j} - \log X_{\cdot\cdot}$, the bias terms $b_i$ and $b_j$ absorb the marginal frequencies, and the dot product again encodes PMI. GloVe can be understood as a weighted least-squares factorization of the log co-occurrence matrix, with the biases playing the role of chemical potentials in the statistical-mechanical analogy.

\subsection{Arora et al.\ (2016): the RAND-WALK model}

Arora et al.~\cite{arora2016} proposed a generative model in which a latent \emph{discourse vector} $\bm{c}_t$ performs a random walk on the unit sphere $S^{d-1}$, and each word $w$ is emitted with log-linear probability
\begin{equation}\label{eq:arora-emission}
    \Pr[w \mid \bm{c}_t] = \frac{1}{Z_{\bm{c}_t}} \exp(\bm{c}_t \cdot \vw_w)\,.
\end{equation}
Under an isotropy assumption on the word vectors, the partition function $Z_{\bm{c}_t}$ concentrates around a constant $\hat{Z}$, and the co-occurrence probability within a window satisfies
\begin{equation}\label{eq:arora-main}
    \log p(w, w') = \frac{\|\vw_w + \vw_{w'}\|^2}{2d} - 2\log \hat{Z} \pm \epsilon\,.
\end{equation}
Expanding $\|\vw_w + \vw_{w'}\|^2 = \|\vw_w\|^2 + 2\,\vw_w \cdot \vw_{w'} + \|\vw_{w'}\|^2$ and using $\log p(w) \approx \|\vw_w\|^2 / (2d) - \log \hat{Z}$, one obtains the central result:
\begin{equation}\label{eq:arora-pmi}
    \PMI(w, w') \approx \frac{\vw_w \cdot \vw_{w'}}{d}\,.
\end{equation}
This provides a generative justification---from the random-walk dynamics---for the same PMI$\leftrightarrow$dot-product correspondence that our variational derivation produces from the energy-model side.

\subsection{Hashimoto et al.\ (2016): metric recovery}

Hashimoto, Alvarez-Melis, and Jaakkola~\cite{hashimoto2016} showed that word embeddings can be understood as recovering a latent semantic metric $d(w_i, w_j)$ from co-occurrence data, where co-occurrence probability is a kernel function $K(d(w_i, w_j))$ that decreases with semantic distance---directly analogous to a Boltzmann factor $e^{-\beta\, d^2}$. Under appropriate conditions, the embeddings produced by Word2Vec and GloVe are consistent estimators of this metric structure.

\section{The RAND-WALK Model: Assumptions and Derivation}\label{sec:randwalk}

The variational derivation of Section~\ref{sec:variational} shows that embedding dot products encode PMI from the ``consumer'' side---fitting a softmax model to Boltzmann co-occurrence statistics. Arora et al.~\cite{arora2016} arrive at the same conclusion from the ``producer'' side, via a generative model whose assumptions we now state precisely.

\subsection{Model assumptions}

\paragraph{(A1) Log-linear emission.} Each word $w$ has a time-invariant latent vector $\vw_w \in \mathbb{R}^d$. Given a discourse vector $\bm{c}_t \in \mathbb{R}^d$, the word emitted at time $t$ is drawn from a log-linear distribution:
\begin{equation}\label{eq:rw-emission}
    \Pr[w \text{ emitted at time } t \mid \bm{c}_t] = \frac{1}{Z_{\bm{c}_t}} \exp\!\big(\langle \bm{c}_t, \vw_w \rangle\big)\,,
\end{equation}
where $Z_{\bm{c}_t} = \sum_{w'} \exp(\langle \bm{c}_t, \vw_{w'} \rangle)$ is the partition function. This is a Boltzmann distribution at unit temperature with energy $E(w, \bm{c}_t) = -\langle \bm{c}_t, \vw_w \rangle$.

\paragraph{(A2) Isotropy of word vectors.} The ensemble of word vectors is generated i.i.d.\ as $\vw = s \cdot \hat{\vw}$, where $\hat{\vw} \sim \mathcal{N}(0, I_d)$ is a spherical Gaussian vector and $s$ is a scalar random variable satisfying $\mathbb{E}[s] = \tau = \Theta(1)$ and $s \leq \kappa$ for an absolute constant $\kappa$. The choice $\tau = \Theta(1)$ ensures that $\Pr[w \mid \bm{c}_t]$ is neither too peaked nor too uniform; the dynamic range of word probabilities is roughly $\exp(\kappa^2)$. Isotropy requires $d \ll n$ (the embedding dimension must be much smaller than the vocabulary size), and means the word vectors have no preferred direction: $\frac{1}{n} \sum_w \vw_w \vw_w^\top \approx \frac{\sigma^2}{d}\, I_d$.

\paragraph{(A3) Slow random walk on the discourse vector.} The discourse vector $\bm{c}_t$ performs a random walk on the unit sphere $S^{d-1}$ with:
\begin{itemize}
    \item \textbf{Stationary distribution:} uniform over $S^{d-1}$.
    \item \textbf{Step size:} $\|\bm{c}_{t+1} - \bm{c}_t\|_2 \leq 2\epsilon / \sqrt{d}$ for a small constant $\epsilon$.
\end{itemize}
The slow-walk condition ensures that consecutive discourse vectors are nearly identical, so that $\langle \vw_w, \bm{c}' - \bm{c} \rangle = O(\epsilon^2)$ for any word vector $\vw_w$ of $O(1)$ norm. This models the local coherence of natural language: the topic drifts slowly, and words within a small window share approximately the same discourse context.

\subsection{Partition function concentration}

The technical heart of the RAND-WALK analysis is showing that the partition function $Z_{\bm{c}}$, despite being a sum of heavy-tailed random variables (each $\exp(\langle \vw_w, \bm{c} \rangle)$ is neither bounded nor sub-Gaussian), concentrates tightly around a constant.

\paragraph{(A4) Concentration (Lemma 2.1 of~\cite{arora2016}).} Under the Bayesian prior (A2):
\begin{equation}\label{eq:rw-concentration}
    \Pr_{\bm{c} \sim \mathcal{C}}\!\Big[(1 - \epsilon_z)\, Z \;\leq\; Z_{\bm{c}} \;\leq\; (1 + \epsilon_z)\, Z\Big] \;\geq\; 1 - \delta\,,
\end{equation}
where $\epsilon_z = \tilde{O}(1/\sqrt{n})$ and $\delta = \exp\!\big(-\Omega(\log^2 n)\big)$.

This is \emph{derived} from (A2), not an independent axiom. The key insight is that the mean partition function $\mathbb{E}_{\vw_w}[Z_{\bm{c}}]$ depends on $\bm{c}$ only through $\|\bm{c}\|_2$:
\begin{equation}\label{eq:rw-mean-Z}
    \mathbb{E}_{\vw_w}[Z_{\bm{c}}] = f\!\big(\|\bm{c}\|_2\big) = f(1)\,,
\end{equation}
since $\|\bm{c}\|_2 = 1$ on the sphere. Explicitly, $f(\alpha) = n\,\mathbb{E}_s[\exp(s^2 \alpha / 2)]$, where the expectation is over the norm distribution $s$. Since the mean is exactly the same constant for all $\bm{c} \in S^{d-1}$, concentration of $Z_{\bm{c}}$ around this mean gives (A4). This result provides a theoretical explanation of the empirical phenomenon of \emph{self-normalization} in log-linear language models: treating $Z$ as a constant is justified because it barely varies.

\subsection{Main theorem}

\paragraph{Theorem 2.2~\cite{arora2016}.} Suppose the word vectors satisfy the concentration inequality~\eqref{eq:rw-concentration}, and the window size is $q = 2$. Then:
\begin{align}
    \log p(w, w') &= \frac{\|\vw_w + \vw_{w'}\|^2}{2d} - 2\log Z \pm \epsilon\,, \label{eq:rw-joint} \\
    \log p(w) &= \frac{\|\vw_w\|^2}{2d} - \log Z \pm \epsilon\,, \label{eq:rw-marginal}
\end{align}
where $\epsilon = O(\epsilon_z) + \tilde{O}(1/d) + O(\epsilon^2)$. Together these imply:
\begin{equation}\label{eq:rw-pmi}
    \boxed{\;\PMI(w, w') = \frac{\langle \vw_w,\, \vw_{w'} \rangle}{d} \pm O(\epsilon)\;}
\end{equation}

Since word vectors have $\ell_2$ norm of order $\sqrt{d}$, the leading term $\|\vw_w + \vw_{w'}\|^2 / (2d)$ is $\Theta(1)$, so the error $\epsilon$ is small in comparison.

\paragraph{Proof mechanism.} The derivation proceeds in three steps:
\begin{enumerate}
    \item \textbf{Integrate out the latent variables.} Write $p(w, w') = \mathbb{E}_{\bm{c}, \bm{c}'}[\Pr(w \mid \bm{c})\,\Pr(w' \mid \bm{c}')]$ and use~\eqref{eq:rw-concentration} to replace $Z_{\bm{c}}$ and $Z_{\bm{c}'}$ by the constant $Z$ with multiplicative error $1 \pm O(\epsilon_z)$.
    \item \textbf{Replace $\bm{c}'$ by $\bm{c}$.} The slow-walk assumption (A3) gives $\langle \vw_{w'}, \bm{c}' - \bm{c} \rangle = O(\epsilon^2)$, so $\exp(\langle \vw_{w'}, \bm{c}' \rangle) \approx (1 \pm O(\epsilon^2))\exp(\langle \vw_{w'}, \bm{c} \rangle)$.
    \item \textbf{Evaluate the spherical integral.} The expectation reduces to $\mathbb{E}_{\bm{c} \sim S^{d-1}}[\exp(\langle \vw_w + \vw_{w'}, \bm{c} \rangle)]$. Since $\langle \vw_w + \vw_{w'}, \bm{c} \rangle$ is approximately Gaussian with variance $\|\vw_w + \vw_{w'}\|^2 / d$ for large $d$, the moment generating function gives $\exp(\|\vw_w + \vw_{w'}\|^2 / 2d)$.
\end{enumerate}

\paragraph{General window size.} For co-occurrence within a window of size $q$, the pair $(w, w')$ can appear in any of $\binom{q}{2}$ positions, yielding:
\begin{equation}\label{eq:rw-window}
    \PMI_q(w, w') = \frac{\langle \vw_w,\, \vw_{w'} \rangle}{d} + \log \binom{q}{2} \pm O(\epsilon)\,.
\end{equation}
The additive shift $\log \binom{q}{2}$ is consistent with the empirical observation of Levy and Goldberg~\cite{levy2014} that SGNS solutions satisfy $\PMI(w, w') - \beta = \langle \vw_w, \vw_{w'} \rangle$ for a constant $\beta$ related to the negative sampling rate and window size.

\subsection{Summary of assumptions}

\begin{table}[h]
\centering
\begin{tabular}{@{}clll@{}}
\toprule
\textbf{\#} & \textbf{Assumption} & \textbf{Type} & \textbf{Role} \\
\midrule
A1 & Log-linear emission $\propto \exp(\langle \bm{c}_t, \vw_w \rangle)$ & Axiom & Generative model structure \\
A2 & Isotropy: $\vw = s\hat{\vw}$, $\hat{\vw} \sim \mathcal{N}(0,I)$, $s \leq \kappa$ & Axiom & Partition function concentration \\
A3 & Slow walk: $\|\bm{c}_{t+1} - \bm{c}_t\| \leq 2\epsilon/\sqrt{d}$ & Axiom & Replace $\bm{c}' \approx \bm{c}$ in windows \\
A4 & $Z_{\bm{c}} \approx Z$ for most $\bm{c} \in S^{d-1}$ & Derived & Factor out partition function \\
\bottomrule
\end{tabular}
\caption{Core assumptions of the RAND-WALK model. (A4) is derived from (A2) and high dimensionality; all others are axioms of the generative model.}
\label{tab:assumptions}
\end{table}

\section{Discussion}\label{sec:discussion}

The derivations above converge on a single picture: \textbf{word embeddings arise as solutions to a variational energy-minimization problem, and their inner-product geometry encodes pointwise mutual information.}

The statistical-mechanical perspective makes this natural. The Boltzmann distribution~\eqref{eq:boltzmann} assigns energies to word configurations via the coupling matrix $G_{ij}$. The softmax model~\eqref{eq:proposal} is a mean-field approximation, and minimizing the KL divergence is equivalent to the standard variational free-energy minimization in statistical physics. The embedding vectors are the variational parameters, and their optimal values---determined by the stationary conditions~\eqref{eq:stationary}---encode the interaction structure of the true distribution.

The partition function $\hat{Z}$ plays its usual role: it ensures normalization but complicates computation. The concentration of $\hat{Z}$ in the Arora et al.\ model~\cite{arora2016} is the analogue of self-averaging of the free energy in disordered systems. The shift by $\log k$ in the Levy--Goldberg result~\cite{levy2014} corresponds to adjusting the chemical potential (the reference energy scale) via the number of negative samples.

This perspective suggests that the success of word embedding methods is not accidental but reflects a deep structural fact: natural language co-occurrence statistics are well-described by a pairwise energy model, and the low-rank structure of the PMI matrix (equivalently, the low dimensionality of the embedding space) reflects the fact that semantic similarity is governed by a small number of latent factors---the ``discourse vectors'' of Arora et al., or equivalently, the slow modes of the energy landscape.

\begin{thebibliography}{10}

\bibitem{mikolov2013}
T.~Mikolov, I.~Sutskever, K.~Chen, G.~Corrado, and J.~Dean,
``Distributed Representations of Words and Phrases and their Compositionality,''
\emph{Advances in Neural Information Processing Systems}, vol.~26, 2013.
arXiv:1310.4546.

\bibitem{pennington2014}
J.~Pennington, R.~Socher, and C.~D.~Manning,
``GloVe: Global Vectors for Word Representation,''
\emph{Proceedings of EMNLP}, pp.~1532--1543, 2014.

\bibitem{levy2014}
O.~Levy and Y.~Goldberg,
``Neural Word Embedding as Implicit Matrix Factorization,''
\emph{Advances in Neural Information Processing Systems}, vol.~27, 2014.
arXiv:1405.4053.

\bibitem{levy2015}
O.~Levy, Y.~Goldberg, and I.~Dagan,
``Improving Distributional Similarity with Lessons Learned from Word Embeddings,''
\emph{Transactions of the Association for Computational Linguistics}, vol.~3, pp.~211--225, 2015.
arXiv:1507.01460.

\bibitem{arora2016}
S.~Arora, Y.~Li, Y.~Liang, T.~Ma, and A.~Risteski,
``A Latent Variable Model Approach to PMI-based Word Embeddings,''
\emph{Transactions of the Association for Computational Linguistics}, vol.~4, pp.~385--399, 2016.
arXiv:1502.03520.

\bibitem{arora2017}
S.~Arora, Y.~Liang, and T.~Ma,
``A Simple but Tough-to-Beat Baseline for Sentence Embeddings,''
\emph{Proceedings of ICLR}, 2017.
arXiv:1704.05358.

\bibitem{hashimoto2016}
T.~B.~Hashimoto, D.~Alvarez-Melis, and T.~S.~Jaakkola,
``Word Embeddings as Metric Recovery in Semantic Spaces,''
\emph{Transactions of the Association for Computational Linguistics}, vol.~4, pp.~273--286, 2016.
arXiv:1608.01622.

\end{thebibliography}

\end{document}
